\documentclass[conference]{IEEEtran}
\IEEEoverridecommandlockouts
% The preceding line is only needed to identify funding in the first footnote. If that is unneeded, please comment it out.
\usepackage{cite}
\usepackage{amsmath,amssymb,amsfonts}
\usepackage{algorithmic}
\usepackage{graphicx}
\usepackage{textcomp}
\usepackage{xcolor}
\def\BibTeX{{\rm B\kern-.05em{\sc i\kern-.025em b}\kern-.08em
    T\kern-.1667em\lower.7ex\hbox{E}\kern-.125emX}}
\begin{document}

\title{Your Cool Project Title}

\author{
\IEEEauthorblockN{Name Surname}
\IEEEauthorblockA{email address}
\and
\IEEEauthorblockN{Name Surname}
\IEEEauthorblockA{email address}
}

\maketitle

NOTE: You have a 4-page limit, not including references and ``LLM Usage" section. (5-page if you are taking the graduate version of the class).

\section{Introduction}
Briefly introduce your project.
\begin{itemize}
    \item Start by setting the context of your problem and the problem statement.
    \item Explain your contributions towards solving the problem.
    \item Summarize three key insights/innovations.
\end{itemize}


\section{Related Works (Graduate Only)}


\section{Your Cool Project}
Describe in some detail how you are solving the problem you have introduced.


\section{Experiment Setup}
Describe in some detail how you set up the experiments to test the validity of your claims.

\section{Results and Discussion}
Present your results, analysis, and exploration. This is a big portion of your grade, so here are some tips on organization, what to include, and the expected depth of analysis:
\begin{itemize}
    \item Start by summarizing your results, observations, and conclusions at a high-level.
    \item If it makes sense, organize results into subsections.
    \item For each result (usually in the form of plots or tables):
    \begin{itemize}
        \item introduce the results (\emph{e.g.}, ``Fig. 1 shows the energy consumption of our design over a range of sparsity values.");
        \item point out important features for the readers to observe (\emph{e.g.}, ``note that energy consumption decreases as sparsity increases; however, the decrease flattens becomes less pronounced for higher sparsity.");
        \item analyze the results (\emph{e.g.}, ``as sparsity increases, the energy savings from gating and skipping increases; however, the design of the accelerator, specifically this and that, prevents it from leveraging sparsity higher than 0.99. This fact is evident in Fig. 2, which shows the energy breakdown per-component for sparsity value 0.9 and 0.99. Note that ...");
        \item summarize your insights (\emph{e.g.}, ``this result shows that: (1) leveraging sparsity can lead to significant energy savings, but (2) the design of an accelerator should be tailored to the specific range of sparsity of the applications.")
    \end{itemize}
\end{itemize}

\clearpage
NOTE: You must include the ``LLM Usage" and ``References" sections (see below), but they will \emph{not} count towards your page limit.

\section{LLM Usage}
Please include short answers to the following questions.

\begin{itemize}
    \item How did you use the LLM? 
    \item For which parts of the project was an LLM most helpful?
    \item For which parts of the project was an LLM \textit{not} helpful?
\end{itemize}

\begin{thebibliography}{00}
\bibitem{b1} Author names, ``Paper title,'' Name or journal/conference, vol. XX, pp. XX--XX, Month Year.
\end{thebibliography}

\end{document}
